\label{chp:intro}

The quark model predicts the existence of baryons containing two heavy quarks. The doubly heavy baryons constitute an important and promising system to study the structure of physics of a system containing two heavy quarks as well as for understanding the strong interaction at the non-perturbative domain.

Doubly heavy baryons lie on the focus of theoretical and experimental investigations. Experimentally, the first observation of doubly heavy baryon $\Xi_{cc}^+$ was reported by SELEX Collaboration \cite{xi+}.
The existence of its iso doublet partner $\Xi_{cc}^{++}$ state was observed by LHCb Collaboration in the $\Lambda_c^+K^-\pi^+\pi^+$ spectrum, with the mass $m_{\Xi_{cc}^{++}}=3621.40\pm0.72\pm0.14$ $MeV$ \cite{xi++}. The mass of $\Xi_{cc}^+$ has been updated in \cite{xi++,aaij2,aaij3,aaij4}. In the upgraded experiments planned at LHCb with large data samples additional states predicted by the quark model and new decay channels of doubly heavy baryons may be observed.

The tremendous experimental developments triggered theoretical studies on the physics of doubly heavy baryons. Mainly the theoretical studies are concentrated on studying the mass, electromagnetic properties, and weak decays of doubly heavy baryons. These studies are critical to understanding the flavor structure and dynamics of these baryons. The mass of doubly charmed baryons within different approaches has been studied in many works such as constituent quark model \cite{consQM1,consQM2,consQM3,consQM4,consQM5,consQM6,consQM8,consQM9}, Regge phenomenology \cite{regge1,regge2}, QCD sum rules method \cite{qcdsr1,qcdsr2,qcdsr3,qcdsr4,AlievOctet,AlievDecuplet} and relativistic quark model\cite{rqm1}. The magnetic moment is another important quantity in studying the inner structure of hadrons. The magnetic moment of doubly charmed baryons within different approaches such as the  quark model, chiral perturbation theory, bag model, effective quark mass scheme, lattice QCD and harmonic oscillator model are calculated in \cite{transition15}, \cite{transition10}, \cite{DhirEMS}, \cite{exBag}, \cite{lattice} and \cite{harmonic} respectively.

The spin-$\nicefrac{3}{2}$ to spin-$\nicefrac{1}{2}$ doubly heavy baryon transition magnetic moments receive special attention as it probes the inner structure as well as possible deformation of heavy baryons.

In the present work, we calculate the transition magnetic moments of doubly heavy baryons with spin-$\nicefrac{3}{2}$ to spin-$\nicefrac{1}{2}$ within the light cone QCD sum rules (LCSR). (More about this method can be found in\cite{lcsr1,lcsr2}.) In this method, the operator product expansion (OPE) is performed over twists of non-local operators.

The work is organized as follows. In \textbf{section II} we derive light cone sum rules for the transition magnetic moments of doubly heavy spin-$\nicefrac{3}{2}$ to spin-$\nicefrac{1}{2}$ transition. In \textbf{section III} we present our numerical results on magnetic moments of considered transitions. In this section, we also present a comparison of our results with predictions of other approaches.